\documentclass[11pt]{article}

% ------------------------------------------------
% Packages
% ------------------------------------------------
\usepackage[margin=1in]{geometry}
\usepackage{amsmath}
\usepackage{xcolor}
\usepackage{setspace}
\usepackage{tcolorbox}
\usepackage{enumitem}
\usepackage{hyperref}
\usepackage{amssymb}

\onehalfspacing

% ------------------------------------------------
% Formatting environments
% ------------------------------------------------

% Referee comment box
\newtcolorbox{commentbox}{
colback=gray!10,
colframe=gray!40,
boxrule=0.5pt,
arc=2pt,
left=6pt,
right=6pt,
top=6pt,
bottom=6pt
}

% Author response
\newenvironment{response}
{\begin{quote}\color{black}}
{\end{quote}}

% Quoted manuscript change
\newtcolorbox{revisionbox}{
colback=white,
colframe=black,
boxrule=0.5pt,
arc=2pt,
left=6pt,
right=6pt,
top=6pt,
bottom=6pt
}

\newenvironment{subcaption2}
{\strut
\vspace{-5pt}
\begin{minipage}[b]{0.95\textwidth}
  \hspace*{-\parindent}
  \footnotesize}
 {\end{minipage}}

\usepackage[backend = biber,
            style=authoryear-icomp,
            citestyle=authoryear-comp,
            url=false,
            isbn=false,
            doi=true,
            % doi=false,
            maxcitenames=2,
            uniquelist=false,
            date=year,
            uniquename=false,
            sorting=ynt,
            sortcites]{biblatex}

\addbibresource{5_Literature/Literature.bib} 

% ------------------------------------------------
% Title
% ------------------------------------------------

\title{Response to the Referees\\
\textbf{JEEM-D-25-01276} \\ 
The heterogeneous effects of climate policy on households: Evidence from 88 countries\footnote{Initially submitted as "Compensation design for carbon pricing with horizontal heterogeneity: Evidence from 88 countries".}
}

\author{Leonard Missbach, Jan Christoph Steckel}

\date{\today}

\begin{document}

\maketitle

\clearpage

% =================================================
% EDITOR
% =================================================

\section*{Response to Referee 1}

\subsection*{General comment}\label{comment_R1_0}

\begin{commentbox}
    This paper is an investigation of the heterogeneous burdens that climate policy imposes on households and, more specifically, the predictors of those household-level burdens. The aim is to use this investigation to inform the design of compensation schemes to offset unequal costs of climate policy. A key novelty of the paper is its geographic scope: the authors collect and analyze data from 88 countries, whereas most of the related existing literature has (I believe) a single-country scope. The key finding advertised by the authors is large heterogeneity, the majority of which is categorized as heterogeneity across the "horizontal" distribution rather than the "vertical distribution"; the latter refers to the distribution of household income quantile, while the former refers to differences within (or conditional on) income quantile. The implication put forward by the authors is that compensation schemes to correct for undesirable distributional effects of climate policies are far from one-size-fits-all and likely do poorly when they fail to consider differences in households beyond income.

    I find it quite believable that income on its own is a poor proxy for climate policy burden, and I think the environmental economics consensus does fail to internalize this, to the detriment of the discipline's recommendations. I also appreciate the ambitious effort to study the question of equity-focused compensation schemes across a huge swath of the world. That said, I think the paper is difficult to read and not streamlined to prove the key point. Moreover, the challenges of integrating data from 88 countries makes it hard for me trust the results related to feature importance and country clustering. I don't expect the authors to reduce the geographic scope of the paper, but I think the paper could be significantly more compelling with a focused effort to clarify the main point and trim away less relevant components.
    Below are a few overarching comments, followed by some more specific ones.
\end{commentbox}

\begin{response}
    Thank you very much for taking time and effort to engage with our work. Your comments were helpful in motivating us to think more carefully about our main messages and methods deployed. We have addressed all of them and think that the resulting changes have contributed to improving the manuscript in multiple dimensions. As a short preview: We have refined our analysis of heterogeneity within and between income groups and put more emphasis on results from our machine learning modeling approach. The clustering exercise, in contrast, now has comparably less prominence and mainly serves a descriptive purpose. We were also paying close attention to connect our results and discussion more closely to our framework, which, as we believe, is important to motivate our research question of interest. 
\end{response}
    
\subsection*{Comment 1}\label{comment_R1_1}

\begin{commentbox}
    Treatment of vertical vs. horizontal distributional effects: The key measure of horizontal heterogeneity in the descriptive analysis is differences in household carbon intensity within expenditure quintile. I do not have command of the literature on vertical vs. horizontal distributional impacts, and I'm willing to believe that "across vs. within quantile" is how the distinction is made. But to me, what matters is the income (or better yet wealth) vs. non-income distinction. Thus the fact that there is more heterogeneity in carbon intensity within quintile than across quintile is misleading, because some of that within-quintile heterogeneity is nonetheless a function of income. Put differently, if a government knows household expenditure/income at a finer level than quintile, then it can use that info to improve its compensation scheme. By the same token, quintile seems arbitrary; what would the key figure(s) look like if decile or something even finer were used?
\end{commentbox}

\begin{response}

We would like to thank you for pointing out important aspects of our analysis that merit further clarification. In response to this comment -- and your related \nameref{comment_R1_E} -- we have decided to substantiate our analysis in the following way:

First, it is correct that most of the literature on distributional effects is primarily concerned with differences between income quintiles, as argued in line 181 and following. Instead, we hope to provide arguments for widening the analytical lens and studying overall heterogeneity instead. 

We now start our descriptive analysis by showing that differences between income groups (expenditure quintiles) are generally insufficient to explain overall heterogeneity. Here, a novel analysis using the Theil index demonstrates that the contribution of variation between expenditure quintiles to overall heterogeneity is generally low (2\%). We corroborate this analysis with three robustness checks: One uses a variant of the Theil index (GE(2)) that specifically accounts for long right-tail distributions; the other two compute and decompose the Theil index by expenditure deciles and expenditure percentiles, specifically. These robustness checks respond to your suggestion to scrutinize our analysis against a finer resolution of expenditure groups. The resulting changes are as follows:

\end{response}

\begin{revisionbox}
    We are interested in overall heterogeneity and to what extent variation in income can help explaining such heterogeneity. The Theil index is a measure for \textit{overall} heterogeneity and can be additively decomposed in two components describing variation between and within groups, respectively \autocite{Shorrocks.1980,Cowell.2011a}. Figure 1 visualizes the Theil index for the carbon intensity of consumption in all countries of our sample. Bars show the share of variation between and within expenditure quintiles, respectively, that sum up to overall heterogeneity. 

% CHANGE: Figure \ref{fig:fig_1} visualizes the distribution of carbon intensity of consumption among the poorest quintile in all countries in our sample. Boxes and whiskers contain 90\% of all households in each quintile and represent horizontal heterogeneity, i.e., differences among poorer households. In contrast, colored bars show the difference between the lowest and highest median carbon intensity across all quintiles for each country, describing vertical heterogeneity, i.e., differences between poorer and richer households\footnote{See Figure \ref{fig:Quint} for country-level comparisons across all expenditure quintiles and Table \ref{tab:A3} for summary statistics on carbon footprints and carbon intensity of consumption.}.

\end{revisionbox}

\begin{figure}[ht!]
    \centering
    %\includegraphics{1_Figures/Figure 1/Figure_1_2017_11B.pdf}
    %\includegraphics[scale=0.7]{1_Figures/Figure 1/Legend Figure 1.pdf}
    \includegraphics[]{1_Figures/Figure 1/Figure_1A_2017_11C.pdf}
    %\caption{Vertical differences and horizontal distribution of carbon intensity within poorest quintiles}
    \caption{Heterogeneity in carbon intensity (Theil index) distinguishing between and within expenditure quintile heterogeneity}
    \label{fig:fig_1}
    \begin{subcaption2}
    Bars show the Theil index of carbon intensity of consumption for 88 countries. A Theil index of zero indicates perfect homogeneity and the index increases with increasing overall heterogeneity. Countries are sorted by overall heterogeneity. For each country, bars indicate the heterogeneity that can be attributed to heterogeneity within expenditure quintiles (horizontal heterogeneity) and to heterogeneity between expenditure quintiles (vertical heterogeneity). See also Table C.9.
    
    % Figure \ref{fig:Quint} shows the distribution of carbon intensities for all expenditure quintiles across all countries in our sample. We rank countries from bottom left to top right by the median carbon intensity in the poorest quintile. Note the different x-axes in the different panels.
    \end{subcaption2}
\end{figure}

\begin{revisionbox}

Figure 1 shows that within-quintile heterogeneity (95\% on average) exceeds between-quintile heterogeneity (5\% on average) by far and in all countries. This result is robust to different robustness checks including an analysis with expenditure deciles or expenditure percentiles (see Table C.9). This highlights the fact that analyses that rely on differences in income to explain differences in the carbon intensity of consumption may be inadequate because they do not account for differences in carbon intensity at similar income levels. Instead, we suggest including household-level characteristics beyond income in such analyses to provide a more nuanced description of which households' consumption is particularly carbon-intensive.
\end{revisionbox}

\begin{response}

This analysis demonstrates clearly that variation between expenditure groups (quintiles, deciles, percentiles) is generally insufficient to explain overall heterogeneity.

Second, we make use of our \textit{sparse model} approach, in which we use total household expenditures as the single feature for predicting household carbon intensity of consumption. As argued in line 615 and following, predictive power is low across countries, which implies that it is difficult to identify a function of income that approximates the resulting carbon intensity of consumption with low error.
\end{response}

\begin{revisionbox}
    We show that variation in total household expenditures alone is often insufficient to predict households' carbon intensity with high accuracy. On average, the goodness of fit (R\textsuperscript{2}) accumulates to 6\% for \textit{sparse} BRT models including only total household expenditures (see Figure B.7 or Table C.12). In 72 countries, such sparse models contribute to explaining no more than 10\% of the variation in carbon intensity. This implies that compensation measures based on household expenditures, such as uniform or differentiated cash transfers, but also reductions in consumption taxes, would prove ineffective in compensating households with the highest additional costs.
\end{revisionbox}

\begin{response}
This approach is consistent with a fine-grained analysis of expenditure groups including non-linear effects, as correctly requested, yet it corroborates our core argument: Variation in household expenditures is a poor predictor for variation in carbon intensity of consumption. In other words, even if a government was to know household expenditures at a very fine resolution, this information will be insufficient to devise an effective compensation measure.
\end{response}

\subsection*{Comment 2}\label{comment_R1_2}

\begin{commentbox}
    Theoretical framework: I can appreciate the implicit pressure in economics to include a theory section in one's paper, but I don't find this six-page section to be worth the space, especially given the technical challenges of the cross-country analysis and the 100-plus page appendix. I think there are some good points about government's challenge of designing a compensation scheme for climate policy, as well as the researcher's challenge of designing a policy-relevant analysis here. But the math seems unnecessary for the rest of the paper, and there are some points that seem irrelevant. For example, do the authors need to assume that government decarbonization targets are exogenous? Is that assumption used anywhere in the rest of the paper. I see that the distributional effects of climate policy can depend on policy stringency, but that seems more like something to acknowledge as a limitation than something to assume away.
\end{commentbox}

\begin{response}

We appreciate your comment, as it addresses an important question that we had asked ourselves when starting to work on this project: How can we substantiate the motivation for our research question? Importantly, we were aiming at pointing out a) why it matters to differentiate vertical and horizontal distributional effects and b) why we should care about identifying the drivers of variation in the first place. This is not trivial, as we believe that our study is more than descriptive. With the help of our model, we argue that identifying household characteristics that predict additional costs matters not only from the pure welfare perspective, i.e., in descriptive terms, but more importantly from the applied policy perspective. One important argument is that each distributional outcome can be achieved with the corresponding set of transfers and we were wishing to demonstrate that such transfers can have different properties in different circumstances -- an argument that has not been sufficiently substantiated. 

In our view, our model includes the core components of the problem and provides the motivation for our subsequent analysis. Nevertheless, we acknowledge that this section was extensive and insufficiently linked to the rest of the analysis. We have streamlined our argumentation where possible. In response to one other comment that we received (Comment 10) -- which seems to appreciate the theoretical model -- we have also ensured to integrate our discussion of results more closely to the presented model.

The assumption of an exogenous climate change mitigation target serves the purpose of restricting the model to its core. Without making this exogenous target explicit, the reader could be tempted to question whether a government should be interested in climate change mitigation at all, given that mitigation incurs heterogeneous effects on the population. In other words, we want to represent that governments cannot respond to reducing distributional effects by not mitigating climate change. (Left aside that unmitigated climate change has distributional effects in itself and a comprehensive representation of the government decision problem would have to factor this in as well.) In the spirit of keeping the theoretical framework both tractable and informative, we have decided to assume the exogenous character of climate change mitigation.

Thank you for suggesting to refine how we discuss policy stringency. By design of our specific analytical approach, the heterogeneity and its drivers are indeed independent of the policy stringency. We have made this more explicit, pointing to our Appendix, where we provide more details:

\end{response}

\begin{revisionbox}
    We derive the heterogeneous costs of climate policy on households by analyzing the heterogeneity in the carbon intensity of consumption: Assume that the consumption of household \textit{A} is twice as carbon-intensive as the consumption of household \textit{B}, then climate policy will lead to twice as high costs for household \textit{A} compared to household \textit{B} and relative to total expenditures\footnote{This proposition holds under the assumption that climate policy raises supply-side input prices in line with the \textit{embedded} CO\textsubscript{2} emissions associated with production, and that firms cannot respond to changing input prices in the short term. As a corollary, output prices for consumer goods and services would increase in proportion  to embedded (direct and indirect) CO\textsubscript{2} emissions. More generally, the carbon intensity reflects the additional cost of any policy that increases consumer prices in proportion to embedded CO\textsubscript{2} emissions, independent of existing policies, but also independent of regulation stringency. One visible example for the equivalence of $e_{i}$ and $c_{i,P^{\ast}}$ is an upstream carbon tax, as we demonstrate in Appendix A.4.}.
\end{revisionbox}

        %\item The assumption of an exogeneous mitigation target are relevant to motivate why the govenment would be interested in mitigation in the first place, i.e., to trim the model.
        %\item As per assumptions, the distributional effects -- as we can analyze them here -- do NOT depend on stringency. Agree that this is a limitation, which, however, is precisely reflected in our assumption.
        %\item Footnote could do the job.
        %\item This comment can be addressed with a mix of explanation and larger (cutting text) and smaller (adding explanations) edits to the text.

\subsection*{Comment 3}\label{comment_R1_3}

\begin{commentbox}
    R-squared improvement from the sparse BRT model to the rich one: If I focus on the average R-squared value reported by the authors in the model accuracy subsection of Section 4, 23\% vs. 4\% is a large relative improvement. But 23\% is not super-inspiring as an absolute measure of predictive power. These numbers really emphasize the point that household expenditure is highly insufficient for the task at hand, but also that including a much broader set of attributes fails to solve the problem. This weakens the contribution of the paper for me.
\end{commentbox}

\begin{response}
    Thank you for your comment which concerns three important aspects of our analysis: First, explanatory power for the \textit{sparse} model is generally low (6 \%), which implies that variation in household expenditures is insufficient to predict variation in carbon intensity of consumption. Second, explanatory power increases substantially (from 6 \% to 24\% on average), when further features or attributes are added for prediction. Third, explanatory power for the \textit{rich} model is still rather low (24 \%).

    We argue that these aspects of our analysis provide important insights for policy design: First, transfers differentiated by income (such as the canonical lump-sum transfers) will be insufficient to target households with the highest additional relative costs -- a finding that, as you write above, is not sufficiently internalized by the environmental economics consensus. Second, addressing heterogeneity -- transcending the vertical dimension of distributional effects -- could benefit from a combination of multiple, differentiated transfers. In some cases, differentiating compensation by factors such as car ownership or cooking fuel use could be substantially more effective than transfers differentiated by income. Third, assuming that governments do not have more detailed consumption data available for analysis, we observe that \textit{any} combination of differentiated transfers will not necessarily be effective in reducing overall heterogeneity. This is because resulting additional costs are more strongly determined by individual household preferences that cannot be observed, such as energy consumption behaviour or preferences for energy efficiency measures. In other words, observable household characteristics are not necessarily a good proxy for household consumption behaviour. 

    In order to minimize heterogeneous effects of climate policy, the government could subsidize the consumption of more carbon-intensive goods and services -- which has by definition a strong predictive power. However, this would counteract the climate policy itself. We have changed our manuscript to reflect such thoughts as follows:

\end{response}

\begin{revisionbox}
    Low model accuracy found for some countries implies that the carbon intensity of consumption cannot be predicted with observable characteristics at hand. This implies that any transfer based on characteristics observable in our dataset will be ineffective in compensating the most carbon-intensive households and in reducing overall heterogeneity. This holds specifically for many countries in Cluster A, underscoring the importance of country- and policy-specific research, especially when governments face information problems \autocite{Mirrlees.1971}.

    In such circumstances, model accuracy would improve drastically by including household-level expenditure shares on different consumption goods and services as predictors. The corresponding compensation policy would consist of reducing excise taxes on comparatively carbon-intensive goods and services. This could contain overall heterogeneity while preserving incentives for supply-side abatement \autocite{Goulder.2008}, even though such a policy would substantially reduce the overall mitigation effect of climate policy. Thus, excise tax cuts may have their merits specifically to address additional market failures, such as technological lock-ins for cooking fuels, preventing substitution to from fossil-based cooking fuels to less expensive, but potentially hazardous biomass \autocite{Greve.2022}.
\end{revisionbox}

\begin{response}
    The bottom line of low model performance in some cases (16 countries lower than 10\% for the \textit{rich} model) is that climate policy will incur substantial heterogeneity, which is difficult to address because of lacking information. We have reflected those thought in the text, e.g., starting in line 630:
    
        % \item Well acknowledged. We discuss this quite openly -- given measurement error and sometimes comparably few data points, a low R\textsuperscript{2} is a feature of the paper and not a bug, as standard economic models imply that R\textsuperscript{2} of household expenditures themselves are large, i.e., transfers differentiated by income are sufficient to address heterogeneity.
        % \item Make this point more explicit in the text and explain it carefully here.
        % \item The referee may like our approach to correct feature importance values with R\textsuperscript{2} for the clustering exercise, which may be, however, shifted to the Appendix.
\end{response}

\begin{revisionbox}
    For some countries, however, the accuracy of rich models is comparatively low. In 16 countries, R\textsuperscript{2} does not exceed 10\%. Model accuracy is lowest in Bulgaria and Suriname (1\%). One reason is that model performance depends critically on data granularity. In cases of low model accuracy, our models are limited to relying on a few available features, such as household expenditures, sub-national area identifiers, household size, or education of the household head. Nevertheless, low model accuracy implies that in some countries it is difficult to infer about households' carbon intensity from observable characteristics, including total household expenditures.
\end{revisionbox}

\begin{response}
    In addition, we would like to highlight our approach for comparing countries to each other in which we normalize feature importance values with R\textsuperscript{2} indicating model performance. This implies that countries will be clustered according to whether features can "explain" variation \textit{overall}, not relative to the available information on household profiles. That is, we believe that low model performance, as observed in some cases, is an important finding of the paper.
\end{response}

\subsection*{Comment 4}\label{comment_R1_4}

\begin{commentbox}
    Variable data availability across countries: The authors are clearly making a concerted effort to optimize given big differences in the scope of available household features across countries. But I am incredulous about the value of the exercise of clustering countries and identifying cross-country feature importance. The overarching takeaway that every country is different is eminently believable. But it's hard for me to believe that the specific results are robust to different availability of household features. It seems likely that governments have access to a more precis and expansive set of household features than the authors in many cases, and it seems quite plausible that if the feature set were to expand for some of the countries with the sparsest sets in this paper, then the results would change meaningfully. I'm sorry that I don't have a specific recommendation for further testing in this vein, but I think more is needed in order to be convincing here.
\end{commentbox}

\begin{response}
    We gratefully acknowledge criticism about the clustering exercise, which is also shared by Reviewer 2 (see their Comment 3). In response, we have substantially shortened this analytical strand, toning down its weight as an analytical finding and label results from the clustering exercise explicitly as descriptive.

    \end{response}

    \begin{revisionbox}
        Country-level analyses can be useful to identify country-specific household characteristics that are associated with higher carbon intensity of consumption. To describe similarities and differences in the importance of characteristics across many countries, we seek to identify clusters of countries. 
    \end{revisionbox}

    \begin{response}
    We were following strict inclusion criteria for the datasets used for this analysis, which led us to exclude many other countries (including China) because important variables were lacking.
    \end{response}

    \begin{revisionbox}
    We include survey datasets in our study if they cover a nationally representative sample, include item-level expenditure information, and if the surveys were conducted between 2010 and 2019\footnote{We exclude more recent survey data, where available, to account for potential biases induced by large economic shocks, such as those associated with Covid-19.}.
    \end{revisionbox}

    \begin{response}
    One challenge is that household expenditure survey data lack important features, which is particularly salient in many European countries because such data is not included in Eurostat's HBS. To address this issue of lacking information in many European countries, we were able to acquire more detailed data for Spain, France, Germany and Romania that include, for example, information about car ownership or heating fuels.  
    
    Nevertheless, we are convinced that our analysis would be interesting, even if it was restricted to the \textit{sparse} model, highlighting that variation in household expenditures is insufficient for explaining variation in carbon intensity of consumption. In general, we can be confident to assume that the data used in our study are the most detailed data that are publicly available for each country including many information about household characteristics that a) we can expect to drive variation in household carbon intensity of consumption or b) are frequently used to devise possible compensation mechanisms for the costs of climate policy, even though -- as we show -- they are insufficient.
    
        % \item Our "overarching takeaway" is an implicit finding that results from our analysis.
        % \item I would doubt that many governments had better data than we do for most of the countries.
        % \item One way forward could be to cut the clustering exercise to a minimum and interpret the results rather than describing in great detail.
\end{response}

\subsection*{Comment 5}\label{comment_R1_5}

\begin{commentbox}
    External validity to the present: I appreciate the authors' acknowledgement that cleaner energy technologies are on the rise, while the data they are using is from 2017 or prior. This limitation is, of course, pervasive in economics research, as retrospective analysis so often uses data from prior to steep rises in renewables penetration and electrification. But this fact (of pervasiveness) doesn't resolve the problem here, and the problem seems particularly large in this context, especially as an interaction with other sources of error that are hard to get around in this global data context.
\end{commentbox}

\begin{response}
    Thank you for this comment to which we respond by making this even more explicit in line 521:
\end{response}

\begin{revisionbox}
    Our approach allows for a consistent, harmonized analysis across countries, but falls short of accounting for the deployment of cleaner technologies since 2017.
\end{revisionbox}

\begin{response}
  You are right that our analysis is somewhat retrospective, but our work includes the most recent consumption data and input-output data available. On the other hand, the industrial structure and trade relationships of countries do change comparably slowly, which would partially mitigate this concern of outdated data.

  In addition, we include the analysis of sectoral policies including an analysis of carbon pricing in the electricity sector, which may serve as a remedy as these findings inform how the deployment of low-emission technologies would change heterogeneity in additional costs of climate policy and in which direction. Accordingly, we have adjusted our manuscript as follows:

        % \item With the respect to external validity to the present, I fear that recently changing consumption patterns have an even greater impact on the analysis compared to changing energy technologies. 
        % \item Discuss role of supplementary analysis with "electricity sector carbon price" as a remedy and illustrate findings more carefully. This analysis can help to discuss how deploying cleaner technologies would change heterogeneity and in which direction.
\end{response}

\begin{revisionbox}
    As a (partial) remedy, we offer analyses for climate policy in the electricity and transport sector, which can help infer about the direction of changes induced by recent technological change.
\end{revisionbox}

\begin{revisionbox}
    For example, policy instruments that lead to marginal price increases for global CO\textsubscript{2} emissions, such as border carbon adjustment \autocite[e.g.,][]{Mehling.2019,Cosbey.2019}, would lead to increasing heterogeneity among richer households relative to poorer households in 60 countries because richer households tend to spend relatively more on imported goods and services. For transport sector policies, we document more carbon-intensive consumption among richer households compared to poorer households in 60 countries, while differences in horizontal heterogeneity exceed vertical differences in 78 countries. 

    In contrast, electricity sector policies are likely to affect poorer households more in 61 countries, with greater horizontal heterogeneity among poorer households in 64 countries. These results also inform about the effects of deploying cleaner technologies. Lowering the CO\textsubscript{2} emissions intensity in the electricity sector is likely leading to less regressive effects and less heterogeneity among poorer households in many contexts. Exemptions include many comparatively poorer countries, where expenditure shares on electricity tend to increase with increasing income. 
\end{revisionbox}

\subsection*{Comment A}\label{comment_R1_A}

\begin{commentbox}
    Input-output data for carbon intensity calculations: My preconceived notion coming into the paper, which I think is not uncommon among environmental economists, is that input-output analysis/data is pretty inaccurate. Can the authors do more to allay this concern?
\end{commentbox}

\begin{response}
    This concern is substantiated given the large aggregation. Yet, it is the most appropriate approach that we are aware in this strand of literature (e.g., \href{https://doi.org/10.1086/701191}{Cronin et al. (2019)}, \href{https://doi.org/10.1016/j.eneco.2021.105550}{Feindt et al. (2021)}, \href{https://doi.org/10.1038/s41893-022-00955-z}{Chancel (2022)}). Our primary goal is to reflect both direct and indirect emissions across industrial sectors and countries, which is comparably difficult to achieve otherwise. Using IO-data allows for a homogenized analysis of inter-sectoral and inter-national production processes. In addition, our choice for the GTAP database is primarily motivated by two aspects: First, the data include a comparably high sectoral resolution (n=65). Second, the data include a comparably high number of low- and middle-income countries (LMIC), which are an important focus of our work. As a result, our work includes analyses of the heterogeneous effects -- beyond vertical heterogeneity -- of climate policy for many countries that had not been covered before. We write:

    \end{response}

\begin{revisionbox}
    We derive total sectoral consumption ($\sum_{i} Y_{i,s,r}$), direct ($E_{s}^{direct}$) and indirect ($E_{s}^{indirect}$) CO\textsubscript{2} emissions from  multi-regional input-output (MRIO) data. This approach is popular among researchers because it accounts for inter-national and inter-sectional differences in production processes and trade flows, while providing sufficient detail for high sectoral resolution. As a result, $e_{s,r}$ expresses CO\textsubscript{2} associated to final household demand for each sector country, accounting for all fossil CO\textsubscript{2} emissions that occur along the entire value chain.
\end{revisionbox}

    % After all, most analyses on the distributional effects of carbon pricing -- which have also laid the foundation to advocate for progressive cash transfers -- use IO-data. Following the same approach and acknowledging that this method has limitations, we arrive at different, important conclusions: Progressive cash transfers are insufficient to alleviate heterogeneous effects of climate policy.
        % \item This concern is substantiated, but standard for those kind of analyses. Luckily, we use IO-data with decent coverage for LMIC, a point that should be strengthened in response.
        % \item Point to other publications in this journal that were relying on IO-data.

\subsection*{Comment B}\label{comment_R1_B}

\begin{commentbox}
    I think the authors should define "horizontal heterogeneity" explicitly in the introduction.
\end{commentbox}

\begin{response}
    Thank you for this suggestion. We agree that that our analysis mandates a greater focus on overall heterogeneity including both vertical \textit{and} horizontal heterogeneity. We have reflected this in the introduction as follows:
    
        % \item I think we should step away from vertical/horizontal or extend this duality by introducing "overall heterogeneity". Introducing these terms in the introduction is an excellent point, let's see, how to integrate it.
\end{response}

\begin{revisionbox}
    In particular, large horizontal heterogeneity \autocite{Fischer.2019,Hansel.2022}, that is, differences in additional costs within income groups, can reduce the effectiveness of uniform lump-sum transfers, which are oft-proposed for their administrative simplicity and progressive effects. [...] We go beyond traditional analyses of vertical and horizontal heterogeneity of policy impacts and use supervised machine learning to disentangle the non-linear contributions of household characteristics to the \textit{overall} variation in the carbon intensity of consumption.
\end{revisionbox}

\subsection*{Comment C}\label{comment_R1_C}

\begin{commentbox}
    On page 3, line 95, the authors define "additional costs $c_i,p <= 0$". Aren't these costs positive? Am I misunderstanding?
\end{commentbox}

\begin{response}
    Indeed, our definition of additional costs may have caused confusion. As you state correctly, the costs are positive. We had chosen to define costs as negative, as costs influence our expression of incidence negatively. We have fixed this by corrections in line 95 and line 147 as follows:
\end{response}
\begin{revisionbox}
    The introduction of climate policy \textit{p} leads to additional costs $c_{i,p}\geq 0$ in household \textit{i}.
\end{revisionbox}

\begin{revisionbox}
The difference between the costs of climate policy and the benefits of compensation policy leads to a na\text{\"i}ve incidence $\pi_{i}$ in household \textit{i}, where

\begin{equation} \label{eq:pi}
    \pi_{i} = \underbrace{-\sum_{p\in P^{\ast}} c_{i,p}}_{\text{costs from climate policy ($\leq 0$)}} + \underbrace{\sum_{t\in T} b_{i,t} + b_{i,\varepsilon}}_{\text{benefits from compensation policy ($\geq 0$)}}
\end{equation}
\end{revisionbox}

\subsection*{Comment D}\label{comment_R1_D}

\begin{commentbox}
    The first full paragraph on page 3 is confusing to me. Pass-through of energy costs is widely documented to vary across space and across industries (see especially Ganapati, Shapiro, and Walker for cross-industry heterogeneity). It sounds like this heterogeneous pass-through is assumed away by the authors. Do I have that right? If so, I'm not sure the authors can avoid such an assumption, but I think there should be a clearer acknowledgement of the issue with that. Pass-through is a component of the first-order approximation of policy impacts on consumer surplus, so it's not nothing to assume it away.
\end{commentbox}

\begin{response}
    Thank you for this comment and correct observation that we do not account for heterogeneous pass-through rates. Acknowledging that our discussion around heterogeneous pass-through rates could have been made more explicit, we now write in line 96 and 882:
\end{response}

\begin{revisionbox}
    This proposition is reasonable for demand-side policies, but also for supply-side policies, assuming that firms pass on additional costs to households\footnote{Manufacturing pass-through rates for changes in energy prices are likely to differ from unity and both between and within industries \autocite{Ganapati.2020,Cludius.2020}. In general equilibrium, such pass-through rates depend inter alia on plant-specific technology use, available substitution technologies, within-industry competition and demand elasticities. In the short-term, adjustments in production technologies may be challenging, such that increases in energy prices can be considered increases in marginal production costs.}.
\end{revisionbox}

\begin{revisionbox}
    Our analysis provides a foundation for more comprehensive analyses using more nuanced data. Such additional research can explicitly address inaccuracies in our modeling approach, including uncertainties about the supply-side pass-through of cost increases, technological path dependencies, and information frictions.
\end{revisionbox}

\begin{response}
    In the short term, it is difficult to see that industries would be able to change production processes and adapt to price signals, especially, if such price signals are comparably small. In the medium-term, pass-through rates depend critically on alternative technologies available as substitution, demand elasticities and sectoral competition -- information which is difficult to obtain on a country-sector-level and about which making assumptions would be incorrect. The impressive work of \textcite{Ganapati.2020}, for example, uses US-American plant-level data between 1972 and 1997 to arrive at its conclusions, but it is unlikely able to inform estimates of country- and sector-specific pass-through rates across our sample. 
    % Discuss more critically the perfect pass-through assumption and explain in the answer that we are thus explicitly studying short-term, first-order figures. I would say that, in the short-term, producers would pass-through any additional cost -- potentially use recent price hikes in oil market as an example (to show that final consumption prices could increase even further than increase in basic input fuels).
\end{response}

\subsection*{Comment E}\label{comment_R1_E}

\begin{commentbox}
    The horizontal distribution coefficient (page 14, line 395) is one sensible measure of differences in expenditure within expenditure quintile. But what about the overall variance of expenditure within quintile? I.e., what if there are long tails but the bulk of households within quintile have very similar expenditures? A measure that captures this seems informative on top of the one chosen by the authors.
\end{commentbox}

\begin{response}

Thank you for this excellent comment! It motivated us to think more carefully about measures of heterogeneity and approaches to clearly distinguish variation with expenditure quintiles from variation between expenditure quintiles. Arguably, and as pointed out correctly, our horizontal distribution coefficient serves a slightly different purpose: It gives an indication of whether horizontal heterogeneity is larger or smaller in the richest quintile compared to the poorest quintile, identifying at what end of the income distribution horizontal heterogeneity matters most.

Importantly, we include a novel analysis in exchange for Figure 1, as shown above:

\end{response}

\begin{revisionbox}
    We are interested in overall heterogeneity and to what extent variation in income can help explaining such heterogeneity. The Theil index is a measure for \textit{overall} heterogeneity and can be additively decomposed in two components describing variation between and within groups, respectively \autocite{Shorrocks.1980,Cowell.2011a}. Figure 1 visualizes the Theil index for the carbon intensity of consumption in all countries of our sample. Bars show the share of variation between and within expenditure quintiles, respectively, that sum up to overall heterogeneity. 

Figure 1 shows that within-quintile heterogeneity (95\% on average) exceeds between-quintile heterogeneity (5\% on average) by far and in all countries. This result is robust to different robustness checks including an analysis with expenditure deciles or expenditure percentiles (see Table C.9). This highlights the fact that analyses that rely on differences in income to explain differences in the carbon intensity of consumption may be inadequate because they do not account for differences in carbon intensity at similar income levels. Instead, we suggest including household-level characteristics beyond income in such analyses to provide a more nuanced description of which households' consumption is particularly carbon-intensive.
\end{revisionbox}

\begin{response}
The former Figure 1 motivated our focus on heterogeneity within expenditure quintiles or heterogeneity that cannot be explained with differences in total household expenditures. Instead, we now calculate the \textbf{Theil index} of carbon intensity of consumption for each country and use its favourable properties to distinguish between heterogeneity that can be attributed to heterogeneity \textit{between} expenditure quintiles and heterogeneity that can be attributed to heterogeneity \textit{within} expenditure quintiles. In addition, this analysis helps identifying countries in which heterogeneity in carbon intensity of consumption is larger than in others (see second column "GE(1)" in Table C.9) -- information that could not be extracted from our previous analysis at first sight.

This analysis corroborates that heterogeneity between expenditure quintiles contributes -- on average -- only 5\% to overall heterogeneity. As a robustness check, we also include this analysis with expenditure deciles and expenditure percentiles -- addressing your substantiated Comment 1 --, where we find that heterogeneity between expenditure quintiles contributes 6\% or 8\% to overall heterogeneity. Furthermore, we analyze a variation of the Theil index (GE(2)) that accounts more explicitly for long right-tail distributions. All robustness checks can be found in Table C.9 and corroborate our core argument that between expenditure group heterogeneity is insufficient to analyze overall heterogeneity.

        % \item Excellent comment, agree, we should think about it and potentially include it.
        % \item Glad that the referee did not dislike this concept altogether.

\end{response}

\subsection*{Comment F}\label{comment_R1_F}

\begin{commentbox}
    Page 14, line 417, the authors write that the BRT algorithm gives "higher weights to observations with larger prediction errors". I would have thought the opposite?
\end{commentbox}

\begin{response}
    Iteratively, the boosted regression tree algorithm fits models on the residuals of the previous model. If the model leads to low prediction error for an observation, the resulting residual for the next iteration is low. If instead the model leads to large prediction error, the resulting residual for the next iteration is large. That is, in each iteration, observations with larger prediction errors are more influential than observations with smaller prediction errors. We acknowledge that the previous formulation has been ambiguous. In response, we have changed the statement in line 425 as follows:
\end{response}

\begin{revisionbox}
    As an extension to regression trees, the BRT algorithm (\texttt{XGBoost} by \textcite{Chen.2016}) fits a sequence of many individual regression trees, iteratively focusing more strongly on observations with larger prediction errors in previous iterations.
\end{revisionbox}

\subsection*{Comment G}\label{comment_R1_G}

\begin{commentbox}
    On the topic of weights: I didn't notice the authors explaining how survey weights are incorporated into the BRT algorithm. In my experience the treatment of weights is super important in machine learning - that is, the results of ML modeling are pretty sensitive to how observations are weighted.
\end{commentbox}

\begin{response}

    Thank you for highlighting the relevance of survey weights. Population surveys use survey weights to ensure that the surveyed sample represents the population. When conducting our analysis using the BRT algorithm, we are in fact interested in the population, not in the sample. Previously and incorrectly, we had not used survey weights for our analysis indeed. In response to your comment, we have analyzed all data including survey weights. We include survey weights for hyperparameter tuning, model fitting and thus the calculation of SHAP values, both for our spare and rich models. Highly skewed survey weights can complicate such operations, so we opt for within-country normalization of survey weights using the average as a denominator. In addition, we trim normalized weights at the 99\textsuperscript{th} percentile.

    All analysis presented in our most recent draft thus include survey weights. In addition, we have added a footnote in the method section as follows:

\end{response}

\begin{revisionbox}
    For each country-level model, we use the full (\textit{rich}) set of household-level characteristics $X_{i}^{'}$ as possible features and perform several feature engineering steps (see also Appendix A.3)\footnote{All models include household survey weights, which we have normalized and trimmed at the 99\textsuperscript{th} percentile level.}.
\end{revisionbox}

\subsection*{Comment H}\label{comment_R1_H}

\begin{commentbox}
    I don't think the authors note what the acronym SHAP (values) stands for.
\end{commentbox}

\begin{response}

We have added the explanation of the acronym as follows:

\end{response}

\begin{revisionbox}
    We also use all observations to evaluate the relative importance of each feature using SHAP (SHapley Additive exPlanations) values \autocite{Lundberg.2017}: Expressed in the unit of the outcome variable, SHAP values represent the contribution of each feature to each individual prediction.
\end{revisionbox}

\subsection*{Comment I}\label{comment_R1_I}

\begin{commentbox}
    The Section 4 title uses the word "determinants" of heterogeneous carbon intensity, which implies causality. The authors do a great job avoiding causal language elsewhere (and the analysis doesn't need causality to be informative), but here I think the word determinants should be changed.
\end{commentbox}

\begin{response}
    That is correct. Using the term "determinants" can be misleading. We have adjusted the title of Section 4, now titling:
\end{response}

\begin{revisionbox}
    Results: Predictors of heterogeneous carbon intensity of consumption
\end{revisionbox}

\begin{response}
    We have also checked carefully for other occurrences of causal language and expressions such as "determinants" and made changes accordingly.
\end{response}
        % \item Well, in this sense, they are in fact determinants of the variation that can be explained. I will screen the respective parts and address it where necessary.

\subsection*{Comment J}\label{comment_R1_J}

\begin{commentbox}
    Footnote 30: The authors write "it can be shown". But it reads to me like it *should* be shown somewhere (in the appendix, though admittedly I don't want the appendix to be longer than it already is), or that a source that does this should be cited.
\end{commentbox}

\begin{response}
    We agree on your suggestion that the Appendix is already quite extensive. However, we are not aware of other work that has demonstrated this argument formally. We are also keen to revise this phrasing, given that this argument does not necessarily need to be shown. Model goodness of fit expresses how much variation in the outcome of interest, i.e., in overall heterogeneity can be explained with various predictors. A set of transfers that would be differentiated according to model outcomes would leave 1-R\textsuperscript{} \% in overall heterogeneity. We thus chose to adjust this footnote as follows, hoping it provides more clarity:
\end{response}

\begin{revisionbox}
    Specifically, goodness of fit (R\textsuperscript{2}) has a convenient interpretation. For any government that chooses a set of compensating transfers, the goodness of fit indicates the maximum possible reduction in overall heterogeneity in \% compared to the policy impact without compensation.
\end{revisionbox}

\begin{response}
    As a brief explanation: R\textsuperscript{2} expresses the variation in carbon intensity of consumption (or the relative additional costs of climate policy) that could potentially be avoided with a set of differentiated transfers. In other words, 1-R\textsuperscript{2} expresses the remaining heterogeneity -- heterogeneity that can be attributed to unobservable household characteristics, such as use of technologies or preferences for more energy-efficient appliances. Thus, it is possible to get an approximation of increasing inequality as a result of climate policy, which could be compared theoretically to the costs for governments to increase their available information on households. It is likely that in many cases the costs of information sourcing does not outweigh the costs of resulting heterogeneity (or resulting changes in public support) -- an application that we would like to propose for future research, yet for a quite narrow audience.
\end{response}    
    % \item It is an elegant mathematical expression, happy to include it in the Appendix or as an answer to this comment, but certainly not prominently in the paper.

\subsection*{Comment K}\label{comment_R1_K}

\begin{commentbox}
    Line 734, page 25, "more compelling" is nebulous and subjective without further explanation.
\end{commentbox}

\begin{response}
    Variation in residential electricity use is strongly correlated with variation in the use of specific appliances. As such information is often unavailable, variation in appliance \textit{ownership} can help infer about variation in electricity use. Indeed, "more compelling" is not sufficiently clear. We have changed the corresponding expression which now reads as follows:
\end{response}  

\begin{revisionbox}
   This is less surprising because appliance ownership is a more relevant, but incomplete, predictor for electricity \textit{use} than electricity \textit{access}.
\end{revisionbox}


\subsection*{Comment L}\label{comment_R1_L}

\begin{commentbox}
    The "identification of country clusters" subsection on page 28 seems like it ought to come before, or at least as Figure 3 is introduced.
\end{commentbox}

\begin{response}
    We have made more explicit that the clustering exercise is primarily performed for descriptive purposes, specifically when we introduce Figure 3.
\end{response}

\begin{revisionbox}
    The importance of features in predicting variation in carbon intensity differs across countries. Figure 3 and Table C.14 show the adjusted feature importance for all features in each country, our vertical distribution coefficient, mean CO\textsubscript{2} intensity, the share of vertical heterogeneity in overall heterogeneity, and R\textsuperscript{2}, grouped by country clusters for illustrative purposes\footnote{See Table C.15 and Figure B.8.2 for the unadjusted feature importance for all features in each country. See Table C.16 and Figure B.8.3 for the adjusted and imputed feature importance for all features in each country.}.
\end{revisionbox}

\begin{response}
    Starting with presenting country-level findings appears to be useful, as the descriptive clustering exercise builds on and synthesizes these findings. The structure of our results section, with a substantially shortened sub-chapter on the clustering results, also follows the structure of our methods section. However, visualizing Figure 3 without clusters, e.g., in alphabetical order would create confusion. Specifically, as we propose to shift Figure 4 to the Appendix, visualizing similarities between countries in this Figure could be informative to the reader.
\end{response}

\subsection*{Comment M}\label{comment_R1_M}

\begin{commentbox}
    Page 33, lines 908-909, I recommend being clearer about what the "equity-pollution dilemma" is.
\end{commentbox}

\begin{response}
    We have changed footnote 42 to provide more information about the equity-pollution-dilemma:
\end{response}

\begin{revisionbox}
    In contrast, research draws attention to the relatively low public acceptance of such transfers and the 'equity-pollution-dilemma'\footnote{The 'equity-pollution-dilemma' describes that reducing inequality, e.g., with progressive transfers could increase emissions on aggregate because the average carbon intensity for poorer households is larger than for richer households \autocite[see also][]{Stiglitz.2019}. In general, redistribution could partially offset demand-side effects of climate policy. Nevertheless, carbon pricing with cost-equivalent transfers to households would be desirable because carbon pricing shifts relative prices between more and less carbon-intensive goods and services.}\autocite{Sager.2019}.
\end{revisionbox}

\subsection*{Comment N}\label{comment_R1_N}

\begin{commentbox}
    Page 34, lines 926-928, I don't follow what "effective" means in the context of labor tax reduction. It sounds like the point is that if the burden falls relatively more on richer households, than lowering labor taxes may offset that. But that defines "effective" in a specific, subjective way, since not everyone wants to offset progressive cost burdens. I might be understanding.
\end{commentbox}

\begin{response}
    We were intending to use the word "effective" to refer to the potential of transfers to address heterogeneity of additional relative costs of climate policy. We understand that this is ambiguous here and have changed the statement as follows:
\end{response}

\begin{revisionbox}
    The discipline has also popularized \textit{reducing distortionary taxes} to reap a 'double dividend' \autocite{Bovenberg.1996}. In addition, lowering income or consumption taxes provides leverage to counter vertical heterogeneity. For example, if richer households consume more carbon-intensively and household expenditures are an important feature, lowering the labor tax may reduce overall heterogeneity effectively.
\end{revisionbox}

\begin{response}
    In this spirit, our discussion of potential compensation mechanisms is generally agnostic of whether progressive effects are desirable or not. We aim to explore options that reduce distortions of any type, even if such distortions may be welfare-improving.
\end{response}
        
        % \item Effective in the sense that it reduces overall heterogeneity of additional impacts. We do not necessarily care about whether effects are regressive or progressive, we just want to minimize overall heterogeneity.
        %\item Can be explained with brief comment, I think.

\subsection*{Comment O}\label{comment_R1_O}

\begin{commentbox}
    Page 34, lines 944-945, don't the authors mean *increase* the price elasticity, not "reduce"? I suppose it depends on whether one is referring to absolute value elasticity or not. But I think we want to encourage better substitutes for the consumption that is being regulated/taxed, which would translate to more price-elastic consumers.
\end{commentbox}

\begin{response}
    Correct, this depends on whether we refer to the absolute value of elasticity or not. We follow your suggestion as follows:
    % \item Depends on how you define it, indeed. Increasing price, increasingly lower demand would be decreasing price elasticity, increasing the absolute value elasticity.
    % \item Make clearer in the text and answer correspondingly.
\end{response}

\begin{revisionbox}
    This can help increase the (absolute) price elasticity of households and make it easier for households to consume less carbon-intensive goods and services.
\end{revisionbox}

\clearpage

\section*{Response to Referee 2}

\subsection*{General comment}\label{comment_R2_0}

\begin{commentbox}
    Thank you for submitting your paper for consideration at JEEM. It was a pleasure reading it. Below I list some recommendations/points for clarification that, in my view, can help strengthen some arguments.
\end{commentbox}

\begin{response}
    Thank you very much for providing constructive feedback, which has helped us to improve our manuscript in clarity and stringency. In the following, we explain how we have addressed your comments in detail. 
\end{response}
    
\subsection*{Comment 1}\label{comment_R2_1}

\begin{commentbox}
    Household-level carbon intensity: My understanding is that you use country-by-sector level data from GTAP to extrapolate household-level emission shares (as a fraction of their total expenditure). Is this correct? In any case, clarifications are needed.
\end{commentbox}

\begin{response}
    Thank you for your collection of comments, which indicate that we were not sufficiently clear on the core of our methodological approach. We address your comments by providing more details in the manuscript about how we calculate direct and indirect effects of carbon intensity. In general, we estimate the household-level carbon intensity of consumption as country-sector-level carbon intensities of final aggregate household demand weighted by household-level consumption shares for different goods and services. Country-sector-level carbon intensities of final household demand comprise both direct carbon emissions of households (e.g., in the transport sector) as well as indirect emissions, i.e., industry emissions that can be attributed to household final demand. This approach ensures that resulting expressions of household carbon intensity of consumption (and thus relative additional costs of climate policy) reflect both changes in resulting fuel prices, but also of price changes for all non-fuel goods and services.

    In the following we reply to your specific points one by one. 

\end{response}

\subsubsection*{Comment 1a}\label{comment_R2_1a}

\begin{commentbox}
    How are emissions from household transportation calculated? I ask about this sector specifically because it seems to be a key driver of heterogeneity (e.g., Figures B.10.2 and B.10.3). Does the process incorporate annualized production emissions, or only fuel use from driving? The only hint that I could get for this question is by looking up the corresponding sector from GTAP, which I believe is: “Land transport and transport via pipelines (otp)” (Appendix, page 61). I wonder if this is more representative of emissions from the transportation of goods (not people).
\end{commentbox}

\begin{response}
    We include both annualized production emissions and fuel use from driving.

    Emissions from household transportation are included in the calculation of household carbon intensity of consumption in multiple ways: First, households report expenditures on transport fuels (such as diesel or gasoline), which we link to the carbon emission intensity of the GTAP sector \textit{petroleum and coke} (\textit{p\_c}). This sector, specifically, captures households direct emissions. 
    
    Second, households report expenditures on transport services (such as public transport, air travel, shipping), which we link to the carbon emission intensity of the respective GTAP sectors \textit{land transport and transport via pipelines} (\textit{otp}), \textit{air transport} (\textit{atp}) and \textit{water transport} (\textit{wtp}), respectively. These sectors capture households' indirect emissions. 
    
    Lastly, households report expenditures on motorized and non-motorized vehicles as assets, which we link to the carbon emission intensity of \textit{manufacture of motorized vehicles} (\textit{mvh}) and \textit{manufacture of other transport equipment} (\textit{otn}), again capturing households' indirect emissions occurring in the manufacturing process.

    We have followed this classification consistently when matching diverse country-specific consumption items to uniform GTAP sectors. Expenditures on public transport fit best to \textit{otp}, which is consistent with substantially lower carbon intensities across countries in comparison to transport fuels (\textit{p\_c}). Importantly, we match single transport fuels at the item-level, i.e., if households would use electricity as a transport fuel, this would be represented in higher expenditures on electricity.

    We hope that this difference between direct and indirect emissions is reflected in sufficient detail in line 290 and following:
    
\end{response}

\begin{revisionbox}
    The carbon intensity of consumption of household \textit{i}, denoted by $e_{i}$, is the variable of interest in this study. It reflects both direct and indirect\footnote{Direct CO\textsubscript{2} emissions refer to emissions resulting from the combustion of fuels in households, e.g., for transportation or heating. Indirect CO\textsubscript{2} emissions refer to emissions that can be attributed to the production, transportation, and retail of all goods and services purchased by each household, such as emissions from electricity generation or manufacturing processes.} CO\textsubscript{2} emissions $E_{i}$ in household \textit{i} relative to total consumption $Y_{i}$ and thus the additional cost of climate policy at the household level $c_{i,p}$, as described in Section 2.
\end{revisionbox}

\subsubsection*{Comment 1b}\label{comment_R2_1b}

\begin{commentbox}
    Are motorcycles and cars assumed to have the same emissions factors?
\end{commentbox}

\begin{response}
    The indirect emissions intensity of motorcycles and cars (as assets) are identical for each country, as we match both consumption items to the same GTAP sector \textit{manufacture of motorized vehicles} \textit{mvh}. This captures emissions occurring in the production process of both motorcycles and cars.

    However, such emissions refer only to expenditures on \textit{buying} motorcycles or cars, but not on \textit{using} them. Heterogeneity between households using a motorcycle or car are reflected in household expenditures on transport fuels. Households driving more inefficient cars would thus report higher expenditures on transport fuels, given the same distance traveled. Importantly, the resulting additional costs of climate policy implicitly represent household use of technology and preference for transport and energy services (see next comment and answer).
\end{response}

\subsubsection*{Comment 1c}\label{comment_R2_1c}

\begin{commentbox}
    Are all cars assumed to have the same emissions factors? This could be problematic for your analysis, in case richer households have more fuel-efficient vehicles, for example.
\end{commentbox}

\begin{response}
    This is an important aspect, but our analysis is in fact agnostic towards individual fuel efficiency of cars. Assume richer households had more fuel-efficient vehicles and they would drive the same distance as poorer households. Then richer households would report lower expenditures on transport fuels, which would in turn lead to a lower household carbon intensity of consumption in comparison to households with the same total expenditures, but less fuel-efficient vehicles. In other words: Our outcome variable of interest is an expression of household use of technology.

    In absence of information on household-specific transport technology (e.g., fuel efficiency of cars) and transport behaviour (e.g., distance driven, individual driving behaviour), our approach captures how much households spend on transport fuels \textit{in comparison} to their total expenditures. This is a relevant variable to estimate the costs that a carbon price (or any other climate policy) would incur on households.
\end{response}

\subsubsection*{Comment 1d}\label{comment_R2_1d}

\begin{commentbox}
    How about the emissions factors assumed for public transport?
\end{commentbox}

\begin{response}
    Expenditures on public transport fits best to \textit{land transport and transport via pipelines} (\textit{otp}), which is consistent with substantially lower carbon intensities across countries in comparison to transport fuels (\textit{p\_c}). We present country-level carbon intensities (of final household demand) in Figure B3.

    It is noteworthy that our calculation of country-sector-level carbon intensities of final demand makes use of information on household final demand originating from GTAP. That is, it accounts for the fact that household final demand for services from this sector refer primarily to {land transport} and not to \textit{transport via pipelines}.
\end{response}

\subsubsection*{Comment 1e}\label{comment_R2_1e}

\begin{commentbox}
    Overall, simplifying assumptions related to the points above might mask substantial heterogeneity across households, so I wonder if this process should be refined with complementary data sources.
\end{commentbox}

\begin{response}
    We are interested in understanding heterogeneity among households in costs incurred by climate policy instruments such as carbon pricing. The carbon intensity of consumption expresses how many kgCO\textsubscript{2} can on average be attributed to one unit of consumption (see Appendix A.4). Our measure reflects whether a household spends money on private transportation (through tracking expenditures on vehicles and carbon-intensive transport fuels) or on public transportation (through tracking expenditures on public transportation), but also how much a household spends on transportation in comparison to all other goods and services. Households which spend relatively more money on transport fuels are more likely to expect higher additional costs in the case of carbon pricing because transport fuels are comparably carbon-intensive (per unit of expenditure). This relationship holds irrespectively of \textit{why} this household spends relatively more money on transport fuels, e.g., because of fuel-inefficient vehicles or because of greater demand for private transportation.

    What is correct is that we attribute carbon emissions and carbon emission intensities based on observed expenditures, but not based on (unobserved) carbon emissions occurring in production. Assume two car producers A and B with the same production process, but B offers their product with a 100\% price increase compared to A. As we -- by definition -- assume homogeneous goods within sectors, this implies that we cannot observe whether a household has bought two vehicles from producer A or one vehicle from producer B. This is a common shortcoming of microsimulations using input-output data. We write:
\end{response}    

\begin{revisionbox}
    Our approach also neglects within-sector differences in the carbon intensity of consumption and relies on consumer price-dependent \textit{expenditures} to calculate household-level carbon intensity instead of the quality and quantity of consumption.
\end{revisionbox}    

    \begin{response}

    Addressing this problem would require plant-level information about production processes and highly disaggregated household consumption data, which we have not come across so far.
    
    Nevertheless, our approach captures differences in production processes between countries, between sectors and differences in household demand for goods and services from different countries and sectors. GTAP includes data for 161 countries and 65 sectors, which strikes a great balance in comparison to other input-output datasets.

    We hope that these reflections have helped to better understand our approach. In response, we have revised the corresponding section on the matching of sector-level carbon intensities of final demand to household-level expenditure shares as follows:
    
\end{response}

\begin{revisionbox}
    We derive total sectoral consumption ($\sum_{i} Y_{i,s,r}$), direct ($E_{s}^{direct}$) and indirect ($E_{s}^{indirect}$) CO\textsubscript{2} emissions from  multi-regional input-output (MRIO) data. This approach is popular among researchers because it accounts for inter-national and inter-sectional differences in production processes and trade flows, while providing sufficient detail for high sectoral resolution. As a result, $e_{s,r}$ expresses CO\textsubscript{2} associated to final household demand for each sector country, accounting for all fossil CO\textsubscript{2} emissions that occur along the entire value chain.
\end{revisionbox}

\begin{revisionbox}
    Our resulting dataset integrates information on household characteristics and household expenditure shares with country- and sector-level carbon intensities, as described in Equation 5. Specifically, it consists of nationally representative accounts of the carbon intensity of household consumption $e_{i}$\footnote{The carbon intensity of household consumption expresses household use of technology and preferences for energy services through sectoral expenditure shares. Everything else equal, switching to cleaner technologies would lead to a lower carbon intensity as household expenditures on fuels would decline. Reducing energy consumption through behavioral change would also result in lower expenditures on fuels, thus lower expenditure shares on fuels and thus a lower carbon intensity of consumption. Thus, without observing household technology and habits, $e_{i}$ represents the relative additional costs of climate policy (see also Appendix A.4), and, most importantly, variation in such costs.}.
\end{revisionbox}

\begin{response}
    As we follow established methodologies for input-output modeling, we feel that this section could in general be kept quite short. Yet, we would be curious to receive further guidance about whether more details about MRIO modeling are required.
\end{response}

\subsection*{Comment 2}\label{comment_R2_2}

\begin{commentbox}
    Machine learning approach: I like the idea of using machine learning to explore heterogeneity and variable importance. However, I have a few recommendations that may help to better leverage advantages of ML.
\end{commentbox}

\begin{response}
    Thank you, that is great to read.
\end{response}

\subsubsection*{Comment 2a}\label{comment_R2_2a}

\begin{commentbox}
    Why did you train country-specific models separately? One of the advantages of ML is flexibility in modeling complex relationships. In theory, you could train a single model that includes (at most 88) country indicators. This should be enough for your tree-based algorithm to “learn” that relationships are different depending on the country. I cannot be sure, but I imagine that this could help improve prediction accuracy for smaller countries (that have a small number of observations for training). I would therefore recommend training a single model, using cross-validation stratified by country (i.e., each CV fold should contain one-fifth of the data from each country).
\end{commentbox}

\begin{response}
    Thank you for this thoughtful comment. Yes, it is possible to train a single model including a country indicator. However, this would assume that there is one common relationship between household characteristic features and the carbon intensity of consumption across countries. Then, using SHAP-based analysis would not be able to identify -- for example -- 88 different curves for household expenditures and the carbon intensity of consumption. We argue that these country-level analyses are precisely what we should be interested in, given that such country-level relationships can inform the design of effective compensation measures. Since compensation for climate policy appears to remain the responsibility of single country governments (the EU may be an exception), we have decided to fit 88 models separately. Estimating a single model, in contrast, would be informative about the design of one uniform, global compensation mechanism across countries.

    We acknowledge that our original manuscript needed improvement regarding the motivation of the analysis and linking Section 2 to the subsequent parts. Throughout the manuscript, we have worked on making this relationship clearer. We hope that these considerations also motivate the estimations of country-level models. 

    In addition and thanks to your suggestion, we have deliberated about the option to include one model for all countries of the EU. Given that our paper is arguably already rich and such an analysis would add yet a new direction while moving away from our core research question, we have decided to leave this approach for future research.

    From a technical perspective, we would also like to highlight that estimating one single model would place more weight on countries with higher populations. This can be critical, given that expenditure levels -- despite homogenization -- may be difficult to compare across countries. All data are crafted with the help of diverse surveys and questionnaires, which implies that data collection may have led to consistent information within countries, but not across countries. Lastly, not all household characteristic features are available in every country, which would limit the model use. These reasons have led us to perform 88 different models, even though some countries may have less observations for training. (\textcite{Elith.2008} suggests that 1.000 observations should be sufficient for meaningful model results. Cambodia, the country with fewest observations in our sample, includes information on 1,207 households.)
\end{response}

\subsubsection*{Comment 2b}\label{comment_R2_2b}

\begin{commentbox}
    I do not see value-added for separate ML models for countries with few available variables or few training observations. For those cases, the advantages of ML are lost, and simpler OLS predictions would perhaps have similar performance.
\end{commentbox}

\begin{response}
    This is a valid concern. Nevertheless and in contrast to OLS, ML models as used in our work can identify non-linear relationships (e.g., between household expenditures and the carbon intensity of consumption) and accommodate for (non-linear) interactions between features. Hypothetical linear relationships would also be captured by BRT. Thus, we have identified benefits that merit the use of BRT, even if there are comparably few features.

    One concern about the use of BRT (in contrast to OLS) can be about overfitting. As a remedy, we perform five-fold cross-validation separately and both for hyperparameter tuning and model fitting. Hyperparameter tuning searches for model parameters out of 100 different combinations. In our previous version, we had been tuning hyperparameters over 16 different combinations. In line 450, we write:
\end{response}

\begin{revisionbox}
    For hyperparameter tuning, we use fivefold cross-validation on each country-level subset of the data; we fit 1,000 trees -- following the recommendations of \textcite{Elith.2008} -- along with 100 different combinations of learning rate, maximum depth of trees, and fraction of features contained in each tree\footnote{We combine different values for learning rate ($\eta \in [0.001,0.1]$), maximum depth of trees (\texttt{max\_depth} $\in \{x \in \mathbb{N} \mid 3  \leq x \leq 15 \}$), and fraction of features contained in each tree (\texttt{mtry} $\in \{0.5,0.7,1\}$). We select randomized combinations of hyperparameters such that the combinations are evenly distributed across the possible combination space using the function \texttt{grid\_space\_filling()} from the \textit{tidymodels} package in R. We show the resulting and preferred combination of hyperparameters in Table ADD.}.
\end{revisionbox}

\subsubsection*{Comment 2c}\label{comment_R2_2c}

\begin{commentbox}
    My recommendation above from point a requires you to restrict the analysis to a set of countries for which the same covariates are available. In my view, this is not an issue, but rather would improve the interpretability of results, putting all countries “on equal grounds.” I would recommend focusing on covariates that are widely available. Also, I wonder if some missing values the “electricity access” variable could be filled by assuming 100\% coverage in developed economies (if this has not been done already).
\end{commentbox}

\begin{response}
    As outlined above in response to your Comment 2a, we are convinced that country-level models are what suits best our analytical angle, given that we perform this analysis primarily to inform the design of (national) compensation mechanisms. And, unfortunately, such compensation mechanisms are unlikely to be implemented at a global, trans-national scale.

    If we understand correctly, your concern specifically addresses the comparison of countries, which may be undue given the lack of some features in some countries. We approach this problem in two ways: First, by normalizing all feature importance variables with overall model accuracy -- this implies that countries are not identified as similar simply because variables are missing. Second, by search for clusters imputing missing values with average corrected feature importance values, leading to comparably few changes.
\end{response}

\begin{revisionbox}
    Based on the adjusted importance of all features and the vertical distribution coefficient of countries, we identify ten distinct clusters of countries\footnote{We present alternative clusters including imputed information for missing values in Figure B.8.3 and Table C.16.}.
\end{revisionbox}

\begin{revisionbox}
    Despite our clustering exercise being comparatively robust to unobserved features (Figure B.8.3), governments may benefit from addressing both hidden information and structurally unobservable predictors that influence households' technology use. 
\end{revisionbox}

\begin{response}
    Your concern about missing data on electricity access is correct. For within-country models, this would however not provide any insights, given that there would be no variation in electricity access. For a cross-country model, this would be useful, again, however, assuming a uniform, cross-country relationship between electricity access and the carbon intensity of consumption, which would be of interest descriptively, but bear little practical significance.

    Please apologize, if our answers on comment 2 appear to be overconfident. After careful reflection, we have decided to advocate for using the national lens as the main angle in this work.
\end{response}

\subsection*{Comment 3}\label{comment_R2_3}

\begin{commentbox}
    Clustering: Why isn’t the horizontal distribution coefficient used for clustering? I think that your framework allows a nice separation of countries into four different cells, depending on vertical and horizontal distribution, as illustrated in Figure 2. Then, different compensation strategies should be considered depending on which cell a given country is part of. It therefore seems natural to focus on clustering along those two dimensions, with other covariates serving only for some refinement.
\end{commentbox}

\begin{response}
    Thank you for this great suggestion. The horizontal distribution coefficient indicates whether horizontal heterogeneity is larger among poorer and richer households and serves mainly a descriptive purpose. Even though two countries appear in the same of four different cells in Figure 2, it does not necessarily follow that effective compensation would need to look similar in both places. 
    
    To better describe heterogeneity between and within expenditure groups -- and in response to two comments from Reviewer 1 -- we have decided to include a more precise measure based on the Theil index and its decomposition. This novel analysis corroborates that heterogeneity between expenditure quintiles is generally insufficient for explaining overall heterogeneity -- an argument that is supported by Figure 2, but not yet sufficiently clear.

    Importantly, we have changed Figure 1, which allows for the following analysis:

\end{response}

\begin{revisionbox}
    We are interested in overall heterogeneity and to what extent variation in income can help explaining such heterogeneity. The Theil index is a measure for \textit{overall} heterogeneity and can be additively decomposed in two components describing variation between and within groups, respectively \autocite{Shorrocks.1980,Cowell.2011a}. Figure 1 visualizes the Theil index for the carbon intensity of consumption in all countries of our sample. Bars show the share of variation between and within expenditure quintiles, respectively, that sum up to overall heterogeneity. 

% CHANGE: Figure \ref{fig:fig_1} visualizes the distribution of carbon intensity of consumption among the poorest quintile in all countries in our sample. Boxes and whiskers contain 90\% of all households in each quintile and represent horizontal heterogeneity, i.e., differences among poorer households. In contrast, colored bars show the difference between the lowest and highest median carbon intensity across all quintiles for each country, describing vertical heterogeneity, i.e., differences between poorer and richer households\footnote{See Figure \ref{fig:Quint} for country-level comparisons across all expenditure quintiles and Table \ref{tab:A3} for summary statistics on carbon footprints and carbon intensity of consumption.}.

\end{revisionbox}

\clearpage

\setcounter{figure}{0}

\begin{figure}[ht!]
    \centering
    %\includegraphics{1_Figures/Figure 1/Figure_1_2017_11B.pdf}
    %\includegraphics[scale=0.7]{1_Figures/Figure 1/Legend Figure 1.pdf}
    \includegraphics[]{1_Figures/Figure 1/Figure_1A_2017_11C.pdf}
    %\caption{Vertical differences and horizontal distribution of carbon intensity within poorest quintiles}
    \caption{Heterogeneity in carbon intensity (Theil index) distinguishing between and within expenditure quintile heterogeneity}
    \begin{subcaption2}
    Bars show the Theil index of carbon intensity of consumption for 88 countries. A Theil index of zero indicates perfect homogeneity and the index increases with increasing overall heterogeneity. Countries are sorted by overall heterogeneity. For each country, bars indicate the heterogeneity that can be attributed to heterogeneity within expenditure quintiles (horizontal heterogeneity) and to heterogeneity between expenditure quintiles (vertical heterogeneity). See also Table C.9.
    
    % Figure \ref{fig:Quint} shows the distribution of carbon intensities for all expenditure quintiles across all countries in our sample. We rank countries from bottom left to top right by the median carbon intensity in the poorest quintile. Note the different x-axes in the different panels.
    \end{subcaption2}
\end{figure}

\begin{revisionbox}

Figure 1 shows that within-quintile heterogeneity (95\% on average) exceeds between-quintile heterogeneity (5\% on average) by far and in all countries. This result is robust to different robustness checks including an analysis with expenditure deciles or expenditure percentiles (see Table C.9). This highlights the fact that analyses that rely on differences in income to explain differences in the carbon intensity of consumption may be inadequate because they do not account for differences in carbon intensity at similar income levels. Instead, we suggest including household-level characteristics beyond income in such analyses to provide a more nuanced description of which households' consumption is particularly carbon-intensive.
\end{revisionbox}

\begin{response}
    We conduct the clustering exercise primarily to indicate options for more effective compensation design. Including measures such as our horizontal distribution coefficient or the Theil-index-based component indicating heterogeneity within expenditure quintiles is not helpful here, as this information is already covered by including the corrected average SHAP values for household expenditures. Countries with high corrected average SHAP values for household expenditures can be more likely to compensate effectively using transfers differentiated by income. Our Theil-index-based component would do not much more than giving higher weight to the variation in the carbon intensity of consumption that can be explained by variation in total household expenditures.

    Including the vertical distribution coefficient, however, is informative, as it indicates whether such transfers would generally need to be progressive (such as the oft-proposed lump-sum transfers) or regressive (e.g., through a labour tax adjustment). In contrast, the corrected average SHAP value of total household expenditures indicates to what extent transfers differentiated by household expenditures can reduce heterogeneous effects altogether. We build on this distinction between total household expenditures as an important feature and direction of relationship between expenditures and carbon intensity in the Discussion.

    In general, including all other features describing household characteristics is central for our approach to identify the most promising options for the design of effective compensation policies.

       % \item I think I have a comment or footnote on this issue somewhere -- possibly, because horizontal heterogeneity is precisely what we wish to explain here. Need to think about it. Potentially, include overall heterogeneity as an additional feature.
       % \item This comment shows, however, that R2 did not fully understand the difference of the descriptive analysis and the clustering exercise. 
       % \item What explains heterogeneity overall (including income - can we use income as an identifier)? Is it progressive or regressive? Theil index: Re-explains contribution of variation in income. Defensible or not. This is why we exclude it, as it would rate again the influence of expenditures on variation. In Figure 4, this does also not need to be reflected, as this information is captured by importance of household expenditures

    Please note that -- in response to comments received by Reviewer 1 -- we have decided to trim the analysis describing country clusters, giving it less weight in the results section. 
\end{response}

\subsection*{Comment 4}\label{comment_R2_4}

\begin{commentbox}
    Title: Since you do not quantify the impact of different compensation designs, I believe that the current title is misleading. Also, I would recommend mentioning both vertical and horizontal heterogeneity in the title.
\end{commentbox}

\begin{response}
    It is true, we do not quantify the impact of any compensation design, even though we could perform this exercise and have done so in a separate web application. 
\end{response}

\begin{revisionbox}
    While we refrain from investigating the importance of household characteristics in predicting the outcomes of such policies for now\footnote{Descriptive analyses for carbon pricing reforms and stylized compensation policies can be accessed and customized through \href{https://www.cpic-global.net/}{a separate webtool}.}, our results highlight that addressing unintended distributional impacts may also have implications for the choice of climate policy instruments, albeit with implications for aggregate efficiency and revenue collection. 
\end{revisionbox}

\begin{response}
    We have decided to change the title of our paper, reflecting our detailed analysis of the overall heterogeneity in the carbon intensity of consumption, transcending vertical \textit{and} horizontal heterogeneity. We now propose:
\end{response}

\begin{revisionbox}
    The heterogeneous effects of climate policy on households: Evidence from 88 countries
\end{revisionbox}

\subsection*{Comment 5}\label{comment_R2_5}

\begin{commentbox}
    Descriptives: The country-level differences from Figure 1 are interesting, but it would be nice to know a few more details about which sectors are mostly driving these patterns. Some descriptives of expenditure shares by country would go a long way (perhaps in the Appendix).
\end{commentbox}

\begin{response}

    Thank you for this attentive suggestion. We have added a novel table that decomposes the country-level carbon intensity of consumption by electricity, transport fuels, other (residential) energy, food, goods and services (Table C.8). We also refer to this table in our manuscript:
\end{response}

\begin{revisionbox}
    Depending on each country, different sectors are more important in driving the carbon intensity of consumption. Table C.8 shows country-level averages and standard deviations for shares of carbon intensities by consumption category, namely electricity, cooking and residential energy, transport fuels, food, goods and services. Across countries, expenditures on services are the most important component of household carbon intensities (23\%), followed by expenditures on transport fuels (22\%) and cooking fuels and other residential energy (21\%). Differences between countries in consumption category shares either reflect country-level differences in expenditure shares on different categories or differing sectoral carbon intensities of production.
\end{revisionbox}

\begin{response}
    In addition, we have put more emphasis on our supplementary analyses that show the heterogeneity in carbon intensity of consumption for sectoral policies in the transport and electricity sector. The refined manuscript parts reads as follows:
    %\item We have a section on these sectors and tons of supplementary material. Will highlight those analyses more carefully and potentially highlight/add further analyses in the Appendix.
\end{response}

\begin{revisionbox}
    For example, policy instruments that lead to marginal price increases for global CO\textsubscript{2} emissions, such as border carbon adjustment \autocite[e.g.,][]{Mehling.2019,Cosbey.2019}, would lead to increasing heterogeneity among richer households relative to poorer households in 60 countries because richer households tend to spend relatively more on imported goods and services. For transport sector policies, we document more carbon-intensive consumption among richer households compared to poorer households in 60 countries, while differences in horizontal heterogeneity exceed vertical differences in 78 countries. 

    In contrast, electricity sector policies are likely to affect poorer households more in 61 countries, with greater horizontal heterogeneity among poorer households in 64 countries. These results also inform about the effects of deploying cleaner technologies. Lowering the CO\textsubscript{2} emissions intensity in the electricity sector is likely leading to less regressive effects and less heterogeneity among poorer households in many contexts. Exemptions include many comparatively poorer countries, where expenditure shares on electricity tend to increase with increasing income. 

    Comparing the coefficients for vertical and horizontal distributional effects across countries and for policy instruments with different regional and sectoral coverage shows that the distributional impacts of climate policy are country- and importantly policy-specific. This implies that governments may prefer some climate policies over others because of distributional concerns, irrespective of the compensation measures available.
\end{revisionbox}

\subsection*{Comment 6}\label{comment_R2_6}

\begin{commentbox}
    Footnote 4: I would move this earlier, around the discussion of assumptions on top of page 5.
\end{commentbox}

\begin{response}
    Agreed. We have moved this footnote around the discussion of assumptions, as suggested. Line 158 now reads as follows:
\end{response}

\begin{revisionbox}
    Households cannot change their household characteristics in the short term, e.g., by improving home insulation or reducing demand for emission-intensive modes of transport\footnote{Available technologies and infrastructure determine the substitution elasticities of households, which, in turn, are a critical determinant of costs to households in the medium-term.}.
\end{revisionbox}

\subsection*{Comment 7}\label{comment_R2_7}

\begin{commentbox}
    Sometimes you refer to “chapters” of the paper, when you should say “sections.”
\end{commentbox}

\begin{response}
    Thank you for pointing out this inaccuracy which has been corrected in all four instances.
\end{response}

\subsection*{Comment 8}\label{comment_R2_8}

\begin{commentbox}
    Assumption on “zero direct emissions from biomass, firewood, and charcoal.” I understand the arguments about informal markets/difficult enforcement for biomass and firewood, but it seems less defensible for charcoal (especially in slightly more developed economies). Can you provide references for this?
\end{commentbox}

\begin{response}
    Unfortunately, we have to rely on this assumption, as the GTAP database includes CO\textsubscript{2} emissions from fossil fuels. CO\textsubscript{2} emissions from the combustion of charcoal, in contrast, are classified as non-fossil CO\textsubscript{2} emissions, in line with the definition by the \href{https://www.ipcc-nggip.iges.or.jp/public/2006gl/pdf/2_Volume2/V2_1_Ch1_Introduction.pdf}{IPCC} (p. 1.15). However, biomass, firewood and charcoal have positive indirect fossil CO\textsubscript{2} emissions, which we account for by matching charcoal to the \textit{lumber} (\textit{lum}) sector. We have added an explanation in the manuscript as follows:
\end{response}

\begin{revisionbox}
    This result is in line with our assumption of zero direct emissions from biomass, firewood, and charcoal because of informal markets and structural impediments to regulating (and taxing) emissions from these sources\footnote{In addition, our CO\textsubscript{2} emissions data only includes emissions from fossil fuels, but not from biomass.}.
\end{revisionbox}

\begin{response}
    Also, in Appendix A.2, we write:
\end{response}

\begin{revisionbox}
    We match items pertaining to fuels such as firewood, charcoal, and other biomass to the sector \textit{lumber} to account for indirect emissions attributable to production, transportation, and retail of these goods. However, we treat direct CO\textsubscript{2} emissions of such fuels as zero, in line with assumptions by the IPCC \autocite{Grad.2023}, but also because direct emissions of such fuels are often difficult to regulate.
\end{revisionbox}
    % \item Will need to look into this, but given that charcoal is not a fossil emissions source, if I am not mistaken, this assumption can be defended following IPCC definitions.

\subsection*{Comment 9}\label{comment_R2_9}

\begin{commentbox}
    Footnote 36: This footnote is longer than a few paragraphs in the main text. Is it really necessary? Try to shorten it, or perhaps move some of this discussion elsewhere.
\end{commentbox}

\begin{response}
    Yes, you are correct. We agree that this footnote is too long and have shortened it in response, putting more weight on the case of Austria:
    %   \item It is, but it is also sort of interesting. I feel that this example illustrates the direct application of our analysis to real world policy design. We can shorten it for sure, though.
\end{response}

\begin{revisionbox}
    With these limitations in mind, we examine which options in the policy toolbox could be more effective in supporting those households that would bear the highest additional costs, thereby reducing horizontal heterogeneity\footnote{Our analysis can shed light on \textit{existing} compensation policies for climate policy. The Austrian government, for example, uses carbon pricing revenues to finance lump-sum transfers, which vary across regions. Households living in regions with less transport and health infrastructure receive higher transfers \autocite{BMK.2023}. Our analysis shows that household expenditures and spatial features account for 52\% of predicted values (R\textsuperscript{2}=0.21). Despite some remaining unobserved heterogeneity and the lack of differentiation of transfers with respect to heating fuels and car ownerships, this transfer mechanism is likely reducing horizontal heterogeneity.}
\end{revisionbox}

\begin{response}
    We believe that this footnote adds information that may be of interest to the reader, which is why we would like to preserve it in general.
\end{response}
    
\subsection*{Comment 10}\label{comment_R2_10}

\begin{commentbox}
    In my view, the current version of Section 5 is disconnected from the nice framework from Section 2. I would recommend making more explicit connections/statements about the impact of different compensation strategies, depending on the current vertical and horizontal distribution coefficients (that are based solely on policy costs).
\end{commentbox}

\begin{response}
    Thank you for this great comment, with which we agree. In response, we have changed Section 5 to connect it more closely to our theoretical framework in Section 2. Examples include:
\end{response}

\begin{revisionbox}
    Here, we refrain from proposing specific compensation policies for particular countries, and recognize that preferences for specific compensation measures over others are subject to normative considerations at the government level. In general, compensation becomes more feasible when climate policy builds on price-based interventions, such as carbon pricing or fossil fuel subsidy removal, thereby increasing the fiscal space to reimburse households. A thorough comparison of different compensation policies should take into account the nature of the existing tax system, the administrative costs of alternative transfer mechanisms, potentially constrained government capacity, and the limited information available to policymakers.
\end{revisionbox}

\begin{revisionbox}
    Following our framework in Section 2, uniform lump-sum transfers would be effective in reducing heterogeneous impacts of climate policy in countries where total household expenditures are an important feature and where disproportionately high costs would fall on relatively poorer households. This is because such transfers would lead to higher benefits for poorer households, thus minimizing the na\text{\"i}ve incidence for all households.
\end{revisionbox}

\begin{revisionbox}
    Uniform cuts in excise taxes on consumption, in contrast, benefit all households equally in relative terms, thus not reducing heterogeneity in the the na\text{\"i}ve incidence at all. Tax reductions differentiated by product (e.g., through VAT), for example through lower tax rates on basic consumer goods, including food and some forms of energy, may reduce heterogeneity, if household expenditures are an important feature and if poorer households consume more carbon-intensively.
\end{revisionbox}

\begin{revisionbox}
    Low model accuracy found for some countries implies that the carbon intensity of consumption cannot be predicted with observable characteristics at hand. This implies that any transfer based on characteristics observable in our dataset will be ineffective in compensating the most carbon-intensive households and in reducing overall heterogeneity of incidence.
\end{revisionbox}

\begin{revisionbox}
    In such circumstances, model accuracy would improve drastically by including household-level expenditure shares on different consumption goods and services as predictors -- as these critically determine the carbon intensity of consumption. The corresponding compensation policy would consist of reducing excise taxes on comparatively carbon-intensive goods and services.
\end{revisionbox}

\begin{revisionbox}
    Addressing the heterogeneous impacts of climate mitigation policy does not necessarily require considering different compensation options. Instead, policymakers can also turn to different types of regulation. [...] While we refrain from investigating the importance of household characteristics in predicting the outcomes of such policies for now\footnote{Descriptive analyses for carbon pricing reforms and stylized compensation policies can be accessed and customized through \href{https://www.cpic-global.net/}{a separate webtool}.}, our results highlight that addressing unintended distributional impacts may also have implications for the choice of climate policy instruments, albeit with implications for aggregate efficiency and revenue collection. 
\end{revisionbox}

\begin{revisionbox}
    The distributional impacts of welfare-enhancing policy proposals are important not only for welfare analyses, but also for understanding the political economy of climate policy. While this study provides a comprehensive assessment of such distributional impacts for climate policy, it is less clear how the distribution of costs translates into public \textit{acceptance}. [...] Future research could contribute to a better understanding of how the (expected) distribution of costs and individual and government-level inequality aversion shapes public acceptance of climate policy. 
\end{revisionbox}

\clearpage

\clearpage

\end{document}